La entrada principal al modelo es la magnitud de la mezcla en dominio STFT, en particular:

\textit{STFT compleja + magnitud + tensor}

Internamente se adapta al formato \textit{[B,T,F]}, se convierte a \textit{dB} y se normaliza antes de entrar a la RNN.
La magnitud lineal se conserva para aplicar la máscara al final.

Los targets ( objetivos ) son las magnitudes de las fuentes reales. Existen dos modalidades de salida:

\begin{itemize}
    \item \textbf{4 fuentes: } \textit{vocales + batería + bajo + resto}
    \item \textbf{2 fuentes: } \textit{vocales , acompañamiento}
\end{itemize}

La salida del modelo no genera audio directamente.Aunque en la ejecución lo parezca, la salida son máscaras tiempo-frecuencia. Posteriormente, se aplica la máscara sobre la magnitud de la mezcla y el resultado corresponde a la estimación de la fuente.
Luego para reconstruir el audio se utiliza la fase de la mezcla. Esto es una limitación importante, ya que es una estimación. La fase de la mezcla original no tiene porque ser la de las fuentes, pero a falta del dato de la magnitud para estas últimas, se utiliza la de la mezcla.


